\chapter{Additional Results}
\label{app:additional-results}

The main text reports compact tables and representative figures. The complete
experiment outputs are stored as repository artifacts so that individual runs,
event-level metrics, and plots can be inspected without expanding the thesis
with very large tables. This appendix documents where those additional results
are located and how they relate to the summaries in \cref{chap:results}.

\section{Native Prediction Metrics}
\label{app:native-prediction-metrics}

Native prediction metrics are stored in the corresponding run folders. These
metrics should be interpreted as diagnostics for the underlying modeling task,
not as replacements for the switch metrics. The main locations are summarized
in \cref{tab:appendix-native-result-locations}.
\Cref{tab:appendix-native-result-locations} is included to make the archived
outputs traceable: for each task it identifies the native metric family and the
repository locations where the full run-level tables can be inspected.

\begin{table}[htbp]
  \centering
  \caption{Repository locations for complete native prediction metrics.}
  \label{tab:appendix-native-result-locations}
  \small
  \setlength{\tabcolsep}{3pt}
  \renewcommand{\arraystretch}{1.15}
  \begin{tabular}{
    @{}
    >{\raggedright\arraybackslash}p{0.22\textwidth}
    >{\raggedright\arraybackslash}p{0.32\textwidth}
    >{\raggedright\arraybackslash}p{0.36\textwidth}
    @{}
  }
    \toprule
    Task & Main metric type & Artifact location \\
    \midrule
    Autoregressive forecasting
      & MAE, RMSE, horizon-wise forecast errors, validation/test summaries.
      & \path{results/runs/autoregressive/} and \path{results/comparisons/model_selection/autoregressive/}. \\
    Current-level persistence
      & Duration-regression MAE, RMSE, median absolute error, and model-selection summaries.
      & \path{results/runs/current_level_persistence/} and \path{results/comparisons/model_selection/current_level_persistence/}. \\
    Long-fade detection
      & AUPRC, AUROC, F1 at the configured probability threshold, and classifier summaries.
      & \path{results/runs/long_fade_detection/} and \path{results/comparisons/model_selection/long_fade_detection/}. \\
    Survival persistence
      & AFT negative log-likelihood, C-index, integrated Brier score, horizon-specific Brier scores, and calibration diagnostics.
      & \path{results/runs/survival_persistence/} and \path{results/comparisons/model_selection/survival_persistence/}. \\
    \bottomrule
  \end{tabular}
\end{table}

For autoregressive models, the saved prediction files contain one row per
forecast window and one column per forecast horizon. This makes it possible to
recompute horizon-wise errors or alternative trajectory-to-switch rules without
retraining the model. For duration and survival models, the prediction files
contain the scalar duration estimate, probability estimate, or survival curve
needed by the downstream switch-conversion scripts.

\section{Switch Metrics and Event-Level Tables}
\label{app:switch-event-tables}

Switch evaluation artifacts are separated from native model metrics. Each
switch-comparison folder contains:

\begin{itemize}
  \item \path{switch_metrics_summary.csv}, with global precision, recall, F1,
  IoU, balanced accuracy, and related binary metrics;
  \item \path{switch_metrics_by_event.csv}, with event-level binary metrics;
  \item \path{global_switch_metrics.csv}, with active duration, number of
  switch islands, and duration agreement diagnostics;
  \item \path{event_switch_metrics.csv}, with the same behavior metrics
  computed per event;
  \item \path{predictions/model_vs_perfect_switch.parquet}, with the aligned
  model and Perfect Switch time series used by the metrics;
  \item \path{figures/event_switch_plots/}, with event-level timelines.
\end{itemize}

\Cref{tab:appendix-switch-artifact-locations} summarizes the folder layout for
these switch artifacts. The distinction between task-native comparisons,
task-level summaries, and cross-task comparisons is important because they use
different timestamp grids and therefore answer different questions.

\begin{table}[htbp]
  \centering
  \caption{Switch-evaluation artifact locations.}
  \label{tab:appendix-switch-artifact-locations}
  \small
  \setlength{\tabcolsep}{3pt}
  \renewcommand{\arraystretch}{1.15}
  \begin{tabular}{
    @{}
    >{\raggedright\arraybackslash}p{0.27\textwidth}
    >{\raggedright\arraybackslash}p{0.65\textwidth}
    @{}
  }
    \toprule
    Comparison type & Artifact location \\
    \midrule
    Task-native model-vs-Perfect comparisons
      & \path{results/comparisons/switch_eval/<task_name>/<selection_id>/<comparison_id>/}. \\
    Task-level model summaries
      & \path{results/comparisons/model_summary/<task_name>/<selection_id>/}. \\
    Cross-task switch comparison
      & \path{results/comparisons/cross_task_switch/<selection_id>/}. \\
    Perfect Switch reference
      & \path{results/switching/perfect_switch/<task_name>/<selection_id>/}. \\
    Central result indexes
      & \path{results/index/}. \\
    \bottomrule
  \end{tabular}
\end{table}

The cross-task comparison uses a common reference grid. Missing native
decisions are mapped to switch \(0\), so the comparison reflects both switch
quality and task coverage. The task-native switch comparisons remain useful
for understanding model behavior on the timestamps where each formulation is
defined.

\section{Additional Diagnostic Plots}
\label{app:additional-diagnostic-plots}

The repository contains diagnostic plots beyond those shown in the main
results chapter. These include forecast examples, worst forecast examples,
horizon-wise error curves, duration-regression scatter plots, event-level
switch timelines, cross-task heatmaps, and native decision-coverage plots.
\Cref{tab:appendix-plot-locations} lists the main plot groups.
\Cref{tab:appendix-plot-locations} is intended as a navigation aid for the
generated visual audit trail. The plots are not all reproduced in the thesis,
but they remain available for checking individual events, model failures, and
coverage differences.

\begin{table}[htbp]
  \centering
  \caption{Main diagnostic plot locations.}
  \label{tab:appendix-plot-locations}
  \small
  \setlength{\tabcolsep}{3pt}
  \renewcommand{\arraystretch}{1.15}
  \begin{tabular}{
    @{}
    >{\raggedright\arraybackslash}p{0.28\textwidth}
    >{\raggedright\arraybackslash}p{0.60\textwidth}
    @{}
  }
    \toprule
    Plot group & Location \\
    \midrule
    Forecast diagnostics
      & \path{results/runs/autoregressive/<selection_id>/<run_id>/figures/}. \\
    Duration-regression diagnostics
      & \path{results/runs/current_level_persistence/<selection_id>/<run_id>/figures/}. \\
    Survival curves and survival diagnostics
      & \path{results/runs/survival_persistence/<selection_id>/<run_id>/figures/}. \\
    Task-native switch timelines
      & \path{results/comparisons/switch_eval/<task_name>/<selection_id>/<comparison_id>/figures/event_switch_plots/}. \\
    Cross-task timelines, heatmaps, and coverage plots
      & \path{results/comparisons/cross_task_switch/<selection_id>/<comparison_id>/figures/}. \\
    \bottomrule
  \end{tabular}
\end{table}

These plots are intentionally kept as generated artifacts rather than copied
exhaustively into the thesis. The main text includes representative examples
that support the empirical discussion, while the artifact folders preserve the
complete visual audit trail for further inspection.
